\chaves{minera\c{c}\~ao de dados educacionais; previs\~ao de desempenho escolar; aprendizado de m\'aquina; risco pedag\'ogico; learning analytics; modelagem temporal tabular.} % Palavras chave separadas por ponto e vírgula

% Não deixe linha em branco após \begin{resumo}
\begin{resumo}
Este Trabalho de Conclus\~ao de Curso investiga como dados gerados pela rotina escolar podem ser organizados para apoiar interven\c{c}\~oes pedag\'ogicas antes que o aluno chegue a um quadro consolidado de baixo desempenho. O problema central consistiu em integrar notas, faltas, pagamentos e informa\c{c}\~oes acad\^emicas em uma pipeline capaz de prever a pr\'oxima nota do aluno em uma disciplina e sinalizar risco pedag\'ogico de curto prazo. Ao longo do projeto, uma hip\'otese inicial baseada em grafos foi testada, mas os experimentos mostraram melhor ader\^encia de uma abordagem temporal tabular, mais compat\'ivel com a natureza longitudinal dos registros escolares e com o tipo de decis\~ao que a escola precisa tomar. A solu\c{c}\~ao final foi implementada no pacote \texttt{school\_predictor}, com separa\c{c}\~ao entre a camada privada de prepara\c{c}\~ao do banco e a camada p\'ublica reproduz\'ivel baseada em arquivos \texttt{CSV}. Sobre essa base, foram realizadas engenharia de atributos, valida\c{c}\~ao temporal por ano letivo, compara\c{c}\~ao entre baselines, Random Forest e HistGradientBoosting, an\'alise de erro e gera\c{c}\~ao de relat\'orios operacionais. Nos experimentos, a previs\~ao num\'erica da pr\'oxima nota obteve melhor resultado com hist\'orico m\'inimo de uma observa\c{c}\~ao anterior por disciplina, alcan\c{c}ando MAE de 0,7659. J\'a o alerta de risco apresentou melhor desempenho com hist\'orico m\'inimo de duas observa\c{c}\~oes, atingindo F1 de 0,7436 e recall de 0,8794 no teste temporal. Como resultado adicional, as sa\'idas anal\'iticas foram convertidas em relat\'orios destinados a professor, coordena\c{c}\~ao e secretaria, incluindo ranking executivo de prioridade. Conclui-se que uma arquitetura temporal voltada ao contexto escolar, quando combinada com interpreta\c{c}\~ao, prioriza\c{c}\~ao e cuidado com dados sens\'iveis, pode ampliar a capacidade institucional de agir de forma mais precoce.
\end{resumo}
